\documentclass[a4paper, serif, 10pt]{moderncv}

%% moderncv
% style and color
\moderncvstyle{banking}
\moderncvcolor{black}

%% page layout/margins
\usepackage[top=2.5cm, bottom=2.5cm, left=1.5cm, right=1.5cm]{geometry}

\nopagenumbers{}

%% Babel config for English language
\usepackage[english]{babel}

%% fontspec
\usepackage{fontspec}

\IfFontExistsTF{Linux Libertine}{
  \setmainfont[Ligatures={TeX, Common}, Numbers=OldStyle]{Linux Libertine}
}{}
\IfFontExistsTF{Linux Libertine O}{
  \setmainfont[Ligatures={TeX, Common}, Numbers=OldStyle]{Linux Libertine O}
}{}

%% CTeX
% CJK support, used to show CJK characters in the resume
%
% - fontset=none: disable builtin fontset but instead set the CJK font manually
% - heading=false: disable ctex heading
% - punct=kaiming: use kaiming punctuations styles for CJK
% - scheme=plain: use plain scheme, do not override \normalsize font size
% - space=auto: space settings for CJK characters
%
% ref:
% - http://ctan.mirrorcatalogs.com/language/chinese/ctex/ctex.pdf
\usepackage[UTF8, heading=false, punct=kaiming, scheme=plain, space=auto]{ctex}

\IfFontExistsTF{Noto Serif CJK SC}{
  \setCJKmainfont{Noto Serif CJK SC}
}{}
\IfFontExistsTF{Noto Sans CJK SC}{
  \setCJKsansfont{Noto Sans CJK SC}
}{}

%% line spacing
\usepackage{setspace}
\setstretch{1.125}

%% URL styling - use normal text instead of monospace
\urlstyle{same}

% Auto-underline all links
% Defer until \begin{document} so hyperref is already loaded
\AtBeginDocument{%
  \let\oldhref\href
  \renewcommand{\href}[2]{\oldhref{#1}{\underline{#2}}}%
}

%% Patch \httplink and \httpslink to handle full URLs
% The original moderncv macros always prepend http:// or https:// to URLs.
% This causes issues when the URL already has a protocol (e.g., https://example.com)
% resulting in malformed URLs like https://https://example.com.
% This patch checks if the URL already contains "://" and uses it directly if so.
% Note: We use \str_set:Nx (x-expansion) to expand macros like \@homepage first.
\ExplSyntaxOn
\str_new:N \g_yamlresume_url_str

\RenewDocumentCommand{\httpslink}{O{}m}{
  \str_set:Nx \g_yamlresume_url_str {#2}
  \str_if_in:NnTF \g_yamlresume_url_str {://}
    { \url{#2} }
    { \url{https://#2} }
}

\RenewDocumentCommand{\httplink}{O{}m}{
  \str_set:Nx \g_yamlresume_url_str {#2}
  \str_if_in:NnTF \g_yamlresume_url_str {://}
    { \url{#2} }
    { \url{http://#2} }
}
\ExplSyntaxOff

\name{Adrian Cole}{}
\title{Data Scientist Specializing in Machine Learning and Predictive Analytics}
\phone[mobile]{(415) 555-0142}
\email{adrian.cole@example.com}
\homepage{https://adriancole.dev}

\address{456 Mission Street, San Francisco, California, United States, 94110}{}{}

\extrainfo{{\small \faGithub}\ \href{https://github.com/adriancole}{@adriancole} {} {} {} • {} {} {}
{\small \faLinkedin}\ \href{https://www.linkedin.com/in/adriancole-data}{@adriancole-data}}

\begin{document}

\maketitle

\section{Basics}

\cvline{}{Data scientist with 6+ years of experience building machine learning systems that turn messy data into measurable business impact. Strong background in statistics, causal inference, and natural language processing, with deep expertise across the full model lifecycle—from problem framing and data pipelines to deployment and monitoring. Passionate about creating interpretable, reliable models and translating data insights into clear recommendations for product and leadership teams.}
\section{Education}

\cventry{Aug 2016–May 2018}
        {Master, Information and Data Science, Score: 3.9}
        {University of California, Berkeley}
        {\url{https://www.ischool.berkeley.edu/}}
        {}
        {\begin{itemize}
\item Concentrated on machine learning, causal inference, and large-scale data systems
\item Capstone project: built a real-time churn prediction pipeline for a telecommunications partner
\item Graduate research assistant supporting a study on algorithmic fairness in lending decisions
\end{itemize}
\textbf{Courses}: Applied Machine Learning, Statistics for Data Science, Natural Language Processing, Big Data Systems, Experimentation and Causal Inference, Data Visualization}

\cventry{Sep 2012–Jun 2016}
        {Bachelor, Statistics, Score: 3.8}
        {University of Washington}
        {\url{https://www.stat.washington.edu/}}
        {}
        {\begin{itemize}
\item Built strong foundations in statistical theory and computational methods
\item Undergraduate research assistant in the Statistical Learning Lab
\end{itemize}
\textbf{Courses}: Statistical Inference, Regression Models, Probability Theory, Computational Statistics, Data Mining}
\section{Work}

\cventry{Mar 2021–Present}
        {Senior Data Scientist}
        {Lumina Analytics}
        {\url{https://www.lumina-analytics.com/}}
        {}
        {\begin{itemize}
\item Lead development of predictive models for customer lifetime value across a portfolio of SaaS clients, improving average revenue forecast accuracy by 18\%
\item Architect feature store and model monitoring pipelines using Python and MLflow, reducing model retraining time from days to hours
\item Partner with product and engineering teams to design A/B tests and causal inference studies for new feature adoption initiatives
\item Mentor a team of four data scientists and establish best practices for code review, reproducibility, and documentation
\end{itemize}
\textbf{Keywords}: Python, MLflow, Feature Store, Causal Inference, Mentorship}

\cventry{Jul 2018–Feb 2021}
        {Data Scientist}
        {HealthInsight}
        {\url{https://www.healthinsight.com/}}
        {}
        {\begin{itemize}
\item Built risk stratification models using electronic health records and claims data, enabling early intervention for high-risk patients
\item Designed and implemented an NLP pipeline to extract clinical concepts from unstructured notes, achieving 87\% F1 score
\item Created interactive Tableau dashboards used by clinical leadership to track quality metrics and reduce hospital readmission rates by 12\%
\end{itemize}
\textbf{Keywords}: Healthcare, NLP, Tableau, Electronic Health Records}
\section{Languages}

\cvline{English}{Native or Bilingual Proficiency \hfill \textbf{Keywords}: IELTS 8.5}
\cvline{Spanish}{Full Professional Proficiency \hfill \textbf{Keywords}: DELE C1}
\section{Skills}

\cvline{Machine Learning}{Expert \hfill \textbf{Keywords}: Python, scikit-learn, XGBoost, TensorFlow, PyTorch}
\cvline{Statistical Modeling}{Advanced \hfill \textbf{Keywords}: R, Bayesian Methods, Time Series, Survival Analysis, Experiment Design}
\cvline{Data Engineering}{Advanced \hfill \textbf{Keywords}: SQL, Spark, dbt, Airflow, AWS, Snowflake}
\cvline{Natural Language Processing}{Advanced \hfill \textbf{Keywords}: Hugging Face, spaCy, Transformers, NER, Topic Modeling}
\cvline{Data Visualization}{Intermediate \hfill \textbf{Keywords}: Tableau, matplotlib, Plotly, ggplot2}
\section{Awards}

\cventry{Dec 2023}
        {Lumina Analytics}
        {Outstanding Contribution in Data Science}
        {}
        {}
        {Recognized for leading the development of a company-wide model monitoring platform that reduced unplanned model retraining incidents by 40\%.}
\section{Certificates}

\cventry{Nov 2021}
        {Amazon Web Services}
        {AWS Certified Machine Learning - Specialty}
        {\url{https://aws.amazon.com/certification/certified-machine-learning-specialty/}}
        {}
        {}

\cventry{Apr 2020}
        {Google}
        {TensorFlow Developer Certificate}
        {\url{https://www.tensorflow.org/certificate}}
        {}
        {}
\section{Publications}

\cventry{Aug 2022}
        {A Scalable Framework for Real-Time Customer Churn Prediction}
        {Proceedings of the ACM SIGKDD Conference on Knowledge Discovery and Data Mining}
        {\url{https://doi.org/10.1145/3523227.3547400}}
        {}
        {\begin{itemize}
\item Presented a production-ready churn prediction framework combining gradient-boosted trees with online feature stores
\item Demonstrated a 23\% relative improvement in AUC versus batch baselines across three industrial datasets
\end{itemize}}
\section{Projects}

\cventry{Jan 2023–Jun 2023}
        {An NLP-powered search and question-answering tool for unstructured clinical documents.}
        {DocQuery}
        {\url{https://github.com/adriancole/docquery}}
        {}
        {\begin{itemize}
\item Developed a retrieval-augmented system using sentence embeddings and a fine-tuned transformer model to answer clinician questions from EHR notes
\item Deployed as a REST API with Docker, reducing average answer lookup time from 20 minutes to under 5 seconds
\end{itemize}
\textbf{Keywords}: NLP, Transformers, Q\&A, Docker, FastAPI}

\cventry{Aug 2020–Dec 2020}
        {Click-through rate prediction model for online advertising campaigns.}
        {AdScore}
        {\url{https://github.com/adriancole/adscore}}
        {}
        {\begin{itemize}
\item Engineered features from impression logs and built a distributed training pipeline in PySpark
\item Achieved a 15\% lift in offline AUC and enabled real-time bidding through model export to ONNX
\end{itemize}
\textbf{Keywords}: CTR, PySpark, Feature Engineering, ONNX}
\section{Interests}

\cvline{Open Source}{Contributing to scikit-learn, Jupyter, Streamlit}
\cvline{Hiking}{Yosemite, Pacific Crest Trail}
\cvline{Chess}{Tournament Play, Puzzle Solving}
\section{Volunteer}

\cventry{Mar 2022–Dec 2023}
        {Data Science Volunteer}
        {Bay Area Data Corps}
        {\url{https://www.bayareadatacorps.org/}}
        {}
        {\begin{itemize}
\item Partnered with local nonprofits to analyze service utilization data and guide resource allocation decisions
\item Built reusable notebook templates to standardize weekly reporting and make findings accessible to nontechnical staff
\end{itemize}}

\end{document}